\chapter{APÊNDICES}

% ----------------------------------------------------------------
% Apêndice A: Arquitetura experimental
% ----------------------------------------------------------------
\chapter{Arquitetura Experimental do Ambiente de Testes}
\label{apendice:arquitetura}

\begin{figure}[h!]
\centering
\caption{Topologia utilizada para testes com Cluster API em ambiente distribuído.}
\includegraphics[width=120mm]{Imagens/cluster_topologia.png}
\fonte{Elaboração própria.}
\end{figure}

\noindent
<<Descreva brevemente os componentes da arquitetura: local do master node, nós gerenciados, regiões envolvidas, nuvem utilizada (se houver).>>

% ----------------------------------------------------------------
% Apêndice B: Comandos e scripts utilizados
% ----------------------------------------------------------------
\chapter{Comandos e Scripts Utilizados na Implementação}
\label{apendice:scripts}

\begin{lstlisting}[language=bash, caption={Script de criação de cluster com Cluster API}, label={lst:script-cluster}]
clusterctl init --infrastructure aws
...
kubectl apply -f cluster-definition.yaml
\end{lstlisting}

\noindent
<<Inclua aqui os principais comandos usados para provisionar o cluster, unir nós e validar o funcionamento.>>

% ----------------------------------------------------------------
% Apêndice C: Resultados detalhados de testes
% ----------------------------------------------------------------
\chapter{Resultados Detalhados de Testes de Latência e Conectividade}
\label{apendice:resultados}

\begin{table}[ht!]
\centering
\caption{Latência média observada entre o nó de controle e os nós gerenciados.}
\begin{tabular}{|c|c|c|}
\hline
\textbf{Região do Nó} & \textbf{Latência Média (ms)} & \textbf{Desvio Padrão (ms)} \\
\hline
Local (on-premise) & 12 & 1.3 \\
AWS (us-east-1)    & 47 & 3.8 \\
AWS (sa-east-1)    & 26 & 2.1 \\
\hline
\end{tabular}
\fonte{Elaboração própria.}
\end{table}

% ----------------------------------------------------------------
% Apêndice D: Análise comparativa (se necessário)
% ----------------------------------------------------------------
%\chapter{Comparação com Ferramentas Alternativas}
%\label{apendice:comparativo}

%\begin{table}[ht!]
%\centering
%\caption{Comparativo entre Cluster API e outras ferramentas de orquestração.}
%\begin{tabular}{|c|c|c|c|}
%\hline
%\textbf{Ferramenta} & \textbf{Orquestração Multi-Cluster} & \textbf{Nível de Abstração} & \textbf{Suporte Industrial} \\
%\hline
%Cluster API & Sim & Alto (Kubernetes-native) & Médio \\
%Rancher     & Sim & Médio                    & Alto \\
%KubeFed     & Sim & Baixo                    & Baixo \\
%\hline
%\end{tabular}
%\fonte{Elaboração própria com base na documentação oficial.}
%\end{table}

% ----------------------------------------------------------------
% Apêndice E: Glossário (opcional)
% ----------------------------------------------------------------
%\chapter{Glossário de Termos Técnicos}
%\begin{itemize}
%    \item \textbf{Cluster API} – Ferramenta para orquestração declarativa de clusters Kubernetes.
%    \item \textbf{Node} – Unidade computacional de um cluster Kubernetes.
%    \item \textbf{Provisionamento} – Processo de criar e configurar recursos computacionais automaticamente.
%\end{itemize}
